\documentclass[11pt]{article}
\usepackage[margin=1in]{geometry}
\usepackage{amsmath,amssymb}
\usepackage[hidelinks]{hyperref}

\title{Math 436 Project: Report Outline\\Plants, Caterpillars, and Bees}
\author{}
\date{}

\begin{document}
\maketitle

\section*{Writing Assignments}

\begin{itemize}
    \item Abstract -- Tristen
    \item Background and Motivation -- Bella
    \item Model 1 -- Ashtyn
    \item Model 2 -- Bella
    \item Interventions -- Bella / Tristen (Bella make plots today)
    \item Analysis and Conclusion -- Ashtyn
\end{itemize}

\section*{New Tasks}

\begin{itemize}
    \item The stated $q$ calculation gives $q=-0.000377$ instead of
    $q=0.0041$. To keep $q=0.0041$, change $\gamma$ to
    $5.96\times10^{-5}$. This balances mortality and competition at the
    reference populations of 40 plants and 10 bees.
    \item With the current parameters,
    $\mu/(\varepsilon a)=0.069/[0.40(0.00138)]=125$ plants.
    \item For Model 1, $K=125$ is the bifurcation point. Below 125, the
    plant-only equilibrium is stable. Above 125, it becomes unstable and
    the plant--caterpillar coexistence equilibrium is positive and stable.
    \item For Model 2, $K>125$ also allows a positive plant--caterpillar
    equilibrium and possibly coexistence.
    \item Make the $K$ bifurcation diagrams cover enough of the range above
    125. Something like 0 to 200 would give us more room to see what happens.
    \item We don't need to choose one fixed $K$; we can look at what happens
    as $K$ varies in our bifurcation diagrams.
\end{itemize}

This is a writing plan for the final report. The main story is to start with
herbivory, add pollination, compare how the models behave, and then test which
interventions help the plants. Write the results and conclusions after the
simulations are finished.

\section{Abstract}

Write this last. Use one paragraph to cover:
\begin{itemize}
    \item We model flowering plants, caterpillars, and bees to study herbivory
    and pollination together.
    \item We start with a plant--caterpillar model, then add bees as mutualists.
    \item We compare equilibrium behavior as carrying capacity changes and
    simulate interventions.
    \item State the main findings: what helped plants, what happened to the
    other populations, and any important limitations.
\end{itemize}

\section{Background and Motivation}

Write about:
\begin{itemize}
    \item Why flowering plants matter and how herbivory and pollination
    affect them.
    \item Why it is useful to consider these interactions together.
    \item The milkweed--monarch--bumblebee example, explaining that it
    motivates a simplified model.
    \item Previous work on predator--prey models and plant--pollinator
    interactions. Use relevant sources Tristen found, after checking the
    citations.
    \item How our project builds on these approaches by examining a resource
    shared by a consumer and a mutualist.
\end{itemize}

End with the questions we will investigate:
\begin{quote}
How does adding bees change plant--caterpillar dynamics? Under what conditions
can the populations persist, and how effectively do bee support, reduced
herbivory, and their combination maintain plant abundance and stability?
\end{quote}

\section{Modeling}

\subsection{Stage 1: Plants and Caterpillars}

Introduce $P$, $C$, and time measured in days. Present the equations:
\begin{align*}
    \frac{dP}{dt} &= rP\left(1-\frac{P}{K}\right)-aPC, \\
    \frac{dC}{dt} &= \varepsilon aPC-\mu C.
\end{align*}

Explain the terms and assumptions:
\begin{itemize}
    \item Plants grow logistically.
    \item Caterpillar consumption depends on both populations.
    \item Plant consumption supports caterpillar population growth.
    \item Caterpillars leave the population at a constant per-capita rate.
    \item Parameters are constant, and spatial differences and life stages
    are omitted.
\end{itemize}

Briefly give the three equilibria and their stability results. Explain that
the proposed shared parameter values give
\[
    \frac{\mu}{\varepsilon a}
    =\frac{0.069}{0.40(0.00138)}=125.
\]

\paragraph{Insert Figure 1: Ashtyn's Model 1 results.}
Use her two $K$-bifurcation panels, updated to the shared parameters and a
range extending beyond 125. Include a representative time simulation as a
small additional panel.

Discuss:
\begin{itemize}
    \item Below 125, the plant-only equilibrium is stable.
    \item Above 125, caterpillars can persist and the plant-only equilibrium
    becomes unstable.
    \item Once caterpillars persist, increasing $K$ supports more caterpillars
    while equilibrium plant abundance remains at 125.
\end{itemize}

Then explain why we introduce Stage 2:
\begin{quote}
Model 1 provides a baseline for herbivory, but cannot answer how pollinators
affect plants or whether supporting bees changes the outcome.
\end{quote}

\subsection{Stage 2: Adding Bees}

Introduce $B$ and present the equations:
\begin{align*}
    \frac{dP}{dt}
        &= rP\left(1-\frac{P}{K}\right)-aPC
        +\beta\frac{vPB}{H+P}\left(1-\frac{P}{K}\right), \\
    \frac{dC}{dt} &= \varepsilon aPC-\mu C, \\
    \frac{dB}{dt} &= \gamma\frac{vPB}{H+P}-dB-qB^2.
\end{align*}

Explain the additional assumptions:
\begin{itemize}
    \item Visitation per bee increases with plant abundance but saturates.
    \item Flower visits benefit plant reproduction and bee population growth.
    \item Plant growth remains limited by $K$.
    \item Bees experience natural mortality and competition.
    \item Bees and caterpillars interact indirectly through plants.
\end{itemize}

Include one parameter table with values, units, and sources or assumptions.
Identify which values are measured estimates and which are calibrated.
Explain that $K$ varies between scenarios while the other parameters remain
fixed during its sweep.

Include the compact equilibrium, existence, and stability table. Briefly
explain how Jacobian eigenvalues determine stability. Keep the entry-by-entry
algebra in the working notes.

\subsection{Numerical Experiments}

Specify the solver, initial conditions, time horizon, and parameters used.
Explain that we will:
\begin{itemize}
    \item Compare Models 1 and 2 using the same shared parameters.
    \item Vary $K$ from a small positive value to around 200.
    \item Compare four intervention scenarios from the same starting state
    and intervention time.
    \item Represent ongoing bee support by reducing $d$, and reduced
    herbivory by lowering $a$, if these remain the chosen interventions.
    \item State the intervention strengths clearly.
\end{itemize}

\section{Results}

\subsection{Comparison as Carrying Capacity Changes}

\paragraph{Insert Figure 2: Model 1 versus Model 2 bifurcation diagrams.}
Use three panels: equilibrium plants, caterpillars, and bees versus $K$.
Distinguish models by color and stable/unstable branches by solid/dashed
lines. Keep $H$, $\gamma$, $q$, and the other parameters fixed while $K$
changes.

Discuss what the actual plots show:
\begin{itemize}
    \item Where each population combination becomes possible.
    \item Whether bees increase the supported caterpillar population.
    \item Whether plant equilibrium abundance changes.
    \item Whether adding mutualism changes stability as $K$ increases.
    \item How Model 2 extends or changes Ashtyn's results.
\end{itemize}

Explain the additional coexistence requirement: the plant abundance
maintained alongside caterpillars must also support bee growth sufficient
to overcome mortality.

\subsection{Interventions to Support Plants}

Choose a representative $K$ from the bifurcation diagram where the
interactions we want to study are active.

\paragraph{Insert Figure 3: Plant abundance over time under four scenarios.}
\begin{enumerate}
    \item No intervention.
    \item Bee support.
    \item Reduced herbivory.
    \item Both together.
\end{enumerate}
Add bee and caterpillar panels if they help explain the plant response.

Discuss:
\begin{itemize}
    \item Which treatment gives higher plant abundance?
    \item Which prevents the lowest plant populations?
    \item Which reduces fluctuations or improves recovery?
    \item Does an initial improvement last?
    \item Do additional caterpillars eventually consume the extra plant
    production?
\end{itemize}

Include a small results table comparing average plant abundance, minimum
plant abundance, and fluctuation size over the same time period.

\subsection{Checks That the Model Behaves Sensibly}

Briefly report checks such as:
\begin{itemize}
    \item With bees initially absent, Model 2 reproduces Model 1.
    \item With both animals absent, plants follow logistic growth.
    \item Numerical equilibrium behavior agrees with the analytical
    predictions.
    \item Populations remain nonnegative.
\end{itemize}
These establish internal consistency. Separately discuss any agreement with
biological evidence.

\section{Analysis and Conclusions}

Answer the original questions using the results. Explain:
\begin{itemize}
    \item What adding bees taught us that Model 1 could not show.
    \item Which intervention helped plants under the conditions tested.
    \item Whether combining interventions provided additional benefits.
    \item Whether benefits depended on environmental capacity or
    intervention strength.
\end{itemize}

Discuss the possibility that more plant production can support more
caterpillars while leaving equilibrium plant abundance unchanged.

Evaluate the assumptions:
\begin{itemize}
    \item One plant variable represents both leaf resources and floral
    resources.
    \item Plants can grow without the modeled bees.
    \item Life stages, seasons, spatial differences, and developmental
    delays are omitted.
    \item Some parameters are rough estimates.
    \item Intervention rankings depend on the chosen model and
    intervention strengths.
\end{itemize}

Describe real revisions we made, including the progression from Model 1 to
Model 2 and any changes prompted by simulations. For example, the bee
calibration needed to account for competition as well as natural mortality.

End with what we would improve with more time: better parameter estimation,
seasonal effects, separate plant biomass or flower resources, or explicit
insect life stages. Keep the conclusions tied to what the simulations
actually establish.

\section{AI Disclosure}

Describe the AI assistance actually used: LaTeX drafting and editing,
equilibrium and stability checks, parameter consistency checks,
simulation planning, and any later coding assistance. Explain what the group
independently checked.

\section*{Final Report Checks}

\begin{itemize}
    \item Use the provided report template without changing its margins or
    fonts. This outline is a planning document, not that template.
    \item Keep the final group report within 10 pages, excluding the
    bibliography. Prioritize the model explanations, compact mathematical
    summaries, and three main figures over the full algebra.
    \item Include a numerical time simulation for each model, not only
    equilibrium diagrams.
    \item Give figures readable labels, units, legends, and captions that
    identify the parameters and explain the main point.
    \item Check that citations match the papers and parameter claims.
    \item Attach a well-documented Jupyter notebook showing the code used.
    \item Leave the heatmap optional until we know whether it would answer
    a question the other figures leave open.
\end{itemize}

\end{document}
